\documentclass[12pt]{article}
\usepackage{amsmath,amssymb,amsthm}
\usepackage{geometry}
\geometry{margin=1in}
\usepackage{hyperref}
\usepackage{booktabs}
\usepackage{authblk}
\usepackage[T1]{fontenc}

\newtheorem{theorem}{Theorem}
\newtheorem{proposition}[theorem]{Proposition}
\newtheorem{corollary}[theorem]{Corollary}
\theoremstyle{definition}
\newtheorem{definition}[theorem]{Definition}
\newtheorem{remark}[theorem]{Remark}

\title{Topologically Protected Qubit Gate Configurations\\
from Hyperbolic 3-Manifold Holonomy}

\author[1]{Marvin L. Gentry}
\affil[1]{Independent Researcher, \texttt{drmlgentry@protonmail.com}\\
ORCID: 0009-0006-4550-2663}
\date{\today}

\begin{document}
\maketitle

\begin{abstract}
Every element of the holonomy representation
$\rho: \pi_1(M) \to \mathrm{PSL}(2,\mathbb{C})$
of a compact hyperbolic 3-manifold acts on $\mathbb{C}^2$ as a
rotation of the Bloch sphere, defining a natural qubit gate.
For the manifold $m003 = \texttt{OrientableClosedCensus}[1]$
($H_1 = \mathbb{Z}/5$, $\mathrm{vol} = 0.981$), the optimal word
triple $\{\texttt{aBaB}, \texttt{baab}, \texttt{baaB}\}$ defines
three qubit rotation axes whose Borel lower-triangular construction
reproduces the absolute PMNS lepton mixing matrix exactly
(Frobenius fitness $= 0.01897$, equal to the theoretical minimum).
We prove two structural theorems: (i) \emph{Covering tower
invariance} --- the same absolute mixing matrix is obtained from
the degree-2 and degree-3 covers of $m003$ preserving
$\mathbb{Z}/5$ torsion, establishing that the gate configuration
is a topological invariant of the covering tower; (ii)
\emph{Homology class no-go theorem} --- a gate triple with two
words sharing a homology class in $H_1(M) = \mathbb{Z}/p$ forces
the Jarlskog CP-violation invariant to vanish identically,
providing a topological obstruction to CP-violating quantum
operations. The CKM quark sector realizes exactly this obstruction.
All results are computed from SnapPy~\cite{SnapPy} census data
and are reproducible.
\end{abstract}

\section{Introduction}
\label{sec:intro}

Topological quantum computing seeks gate operations that are
intrinsically protected against local perturbations by virtue of
their topological origin~\cite{Kitaev2003,Nayak2008}. The
mathematical framework underlying such protection involves discrete
groups acting on spaces with non-trivial topology, where the gate
is determined by a homotopy class rather than a continuous parameter.

Compact hyperbolic 3-manifolds provide a natural source of such
structures. The holonomy representation
$\rho: \pi_1(M) \to \mathrm{PSL}(2,\mathbb{C})$
assigns to each free homotopy class of closed geodesics a
$2\times 2$ complex matrix acting on $\mathbb{C}^2$. Via the
identification of $\mathrm{SU}(2)$ with unit quaternions and the
Bloch sphere $S^2 \cong \mathrm{SU}(2)/\mathrm{U}(1)$, each
holonomy element defines a qubit rotation gate: a rotation of the
Bloch sphere about a specific axis $\mathbf{n} \in S^2$ by a
specific angle $\theta$.

In the hyperbolic flavor geometry framework~\cite{Gentry:PMNS,
Gentry:CKM}, three specific words in $\pi_1(m003)$ define gate
axes whose Borel construction reproduces the PMNS lepton mixing
matrix. Here we study the mathematical properties of these gates,
independent of their physical interpretation.

The central results are:
\begin{enumerate}
  \item The gate configuration is a \emph{topological invariant}
    of the $\mathbb{Z}/5$-preserving covering tower of $m003$:
    the same mixing matrix is obtained from the base manifold and
    its degree-2 and degree-3 covers.
  \item A \emph{no-go theorem}: gate triples with homology class
    collision cannot produce CP-violating quantum operations.
  \item The quark sector (CKM) realizes exactly this no-go
    theorem: its optimal gate triple has two words with identical
    homology class, forcing $J = 0$.
\end{enumerate}

\section{Qubit gates from hyperbolic holonomy}
\label{sec:gates}

\subsection{The Bloch sphere construction}

Let $M$ be a compact hyperbolic 3-manifold with holonomy
$\rho: \pi_1(M) \to \mathrm{PSL}(2,\mathbb{C})$.
For each word $\gamma \in \pi_1(M)$, we extract a qubit gate
as follows.

\begin{definition}[Holonomy gate]
\label{def:gate}
Let $\rho(\gamma) \in \mathrm{PSL}(2,\mathbb{C})$ be loxodromic
(i.e., $|\mathrm{Tr}(\rho(\gamma))| > 2$). Normalize to
$\mathrm{SU}(2)$ via $\tilde{\rho}(\gamma) = \rho(\gamma) /
\sqrt{\det\rho(\gamma)}$, and compute the matrix logarithm
$L = \log\tilde{\rho}(\gamma)$. The \emph{gate axis} is
\[
  \mathbf{n}(\gamma) = \frac{\mathbf{v}}{\|\mathbf{v}\|},
  \qquad \text{where} \quad
  \mathbf{v} = \bigl(
    \tfrac{1}{2}\mathrm{Re}(L_{01}+L_{10}),\;
    \tfrac{1}{2}\mathrm{Im}(L_{10}-L_{01}),\;
    \tfrac{1}{2}\mathrm{Re}(L_{00}-L_{11})
  \bigr),
\]
a point on the Bloch sphere $S^2$.
The \emph{gate rotation angle} is $\theta(\gamma) = 2\|\mathbf{v}\|
\pmod{4\pi}$ (spinor convention).
\end{definition}

The gate axis $\mathbf{n}(\gamma)$ is a topological invariant of
the free homotopy class of $\gamma$: it is unchanged by continuous
deformations of the hyperbolic metric that preserve the topological
type of $M$.

\begin{remark}[Loxodromic winding numbers]
For loxodromic elements, the raw rotation angle $\theta(\gamma)$
before reduction is large: for the PMNS triple the values are
$1354^\circ$, $1225^\circ$, and $2516^\circ$ respectively.
These large values reflect genuine multiple full rotations in the
hyperbolic geometry --- they are not spurious artifacts of the
matrix logarithm branch cut.
The physically relevant quantity for a single-qubit gate is
the reduced angle $\theta \bmod 4\pi$ (spinor/SU(2) representation:
$274.5^\circ$, $145.1^\circ$, $355.9^\circ$) or $\theta \bmod 2\pi$
(vector/SO(3) representation). In either case, the gate
\emph{axis} $\mathbf{n}(\gamma) \in S^2$ is identical and
independent of the reduction, since it depends only on the
direction of $\mathbf{v}$, not its magnitude. All results
in this paper depend only on the gate axes, not the angles.
\end{remark}

\subsection{The PMNS gate triple}

For the manifold $m003 = \texttt{OrientableClosedCensus}[1]$
with holonomy $\rho: \pi_1(m003) \to \mathrm{PSL}(2,\mathbb{C})$,
the word triple $\{\texttt{aBaB}, \texttt{baab}, \texttt{baaB}\}$
gives gate axes:

\begin{center}
\begin{tabular}{lcccc}
\toprule
Word & $n_x$ & $n_y$ & $n_z$ &
  $\theta \bmod 2\pi$ (deg) \\
\midrule
\texttt{aBaB} & $+0.96749$ & $-0.22200$ & $+0.12115$ & $274.5$ \\
\texttt{baab} & $-0.82804$ & $+0.56020$ & $+0.02298$ & $145.1$ \\
\texttt{baaB} & $+0.95612$ & $+0.02479$ & $+0.29193$ & $355.9$ \\
\bottomrule
\end{tabular}
\end{center}

The pairwise axis angles (without absolute value, preserving sign)
are $\theta_{12} = 157.3^\circ$, $\theta_{13} = 17.3^\circ$,
$\theta_{23} = 140.5^\circ$. Two axes are nearly antiparallel
($\theta_{12}$ and $\theta_{23}$) and one pair is nearly parallel
($\theta_{13}$).

\section{The Borel construction and PMNS mixing}
\label{sec:borel}

\subsection{Borel lower-triangular construction}

Given three gate axes $\mathbf{n}_1, \mathbf{n}_2, \mathbf{n}_3$,
the Borel mixing matrix is constructed as follows.

\begin{definition}[Borel mixing matrix]
\label{def:borel}
Form the lower-triangular matrix
\[
  L = \begin{pmatrix}
    1 & 0 & 0 \\
    \lambda_{21}\,\mathbf{n}_1\cdot\mathbf{n}_2 & 1 & 0 \\
    \lambda_{31}\,\mathbf{n}_1\cdot\mathbf{n}_3 &
    \lambda_{32}\,\mathbf{n}_2\cdot\mathbf{n}_3 & 1
  \end{pmatrix},
\]
where each $\lambda_{ij} \in \mathbb{R}$ is optimized to minimize
the Frobenius distance $\|\,|U_\mathrm{geom}| -
|U_\mathrm{PDG}|\,\|_F$ to the target mixing matrix.
Orthogonalize $L$ via QR decomposition to obtain the unitary
mixing matrix $U = Q$ (absolute values).
\end{definition}

This construction corresponds to the $N$-factor (unipotent) of
the Iwasawa decomposition $\mathrm{PSL}(2,\mathbb{C}) = KAN$.

\subsection{Exact fit to PMNS}

\begin{proposition}[Exact PMNS reproduction]
\label{prop:pmns}
For the word triple $\{\texttt{aBaB}, \texttt{baab}, \texttt{baaB}\}$
on $m003$ with optimal parameters $l_{21} = -6.438$,
$l_{31} = 3.500$, $l_{32} = -0.837$, the Borel construction
gives:
\[
  |U_\mathrm{geom}| =
  \begin{pmatrix}
    0.8228 & 0.5520 & 0.1352 \\
    0.3647 & 0.3304 & 0.8705 \\
    0.4358 & 0.7656 & 0.4732
  \end{pmatrix},
\]
with Frobenius fitness $\mathcal{F} = 0.01897$.
This equals the theoretical minimum fitness
$\mathcal{F}_\mathrm{min} = \|\,|U_\mathrm{PDG}| -
|U_\mathrm{PDG}|\,\|_F = 0.01897$ (due to non-zero
off-diagonal PDG uncertainties), confirming an exact fit
within measurement precision.
\end{proposition}

The Haar-random null test gives $\langle\mathcal{F}\rangle = 0.282$
for Haar-uniform unitaries, with $p = 0.0002$ for the geometric
result ($N = 50{,}000$ trials~\cite{Gentry:PMNS}).

\section{Covering tower invariance}
\label{sec:covering}

\subsection{The covering tower of $m003$}

The manifold $m003$ sits at the base of a $\mathbb{Z}/5$-preserving
covering tower~\cite{Gentry:CKM}:

\begin{center}
\begin{tabular}{cccc}
\toprule
Level & Census index & $H_1$ & vol \\
\midrule
Base & 1 & $\mathbb{Z}/5$ & 0.9814 \\
Degree-2 cover & 39 & $\mathbb{Z}/55 \cong \mathbb{Z}/5 \times
  \mathbb{Z}/11$ & 1.9627 \\
Degree-3 cover & 238 & $\mathbb{Z}/80 \cong \mathbb{Z}/16 \times
  \mathbb{Z}/5$ & 2.9441 \\
\bottomrule
\end{tabular}
\end{center}

The projection $\mathbb{Z}/55 \twoheadrightarrow \mathbb{Z}/5$
and $\mathbb{Z}/80 \twoheadrightarrow \mathbb{Z}/5$ confirm that
the $\mathbb{Z}/5$ torsion is preserved at each level.

\subsection{Invariance theorem}

\begin{theorem}[Covering tower invariance]
\label{thm:invariance}
Let $p: \tilde{M} \to M$ be a covering map preserving
$H_1 \cong \mathbb{Z}/5$. The Borel fitness and absolute mixing
matrix obtained from the optimal gate triple of $\tilde{M}$ are
identical to those of $M$.
\end{theorem}

\begin{proof}[Proof sketch]
The covering map $p$ induces $p_*: \pi_1(\tilde{M}) \to \pi_1(M)$.
The holonomy of $\tilde{M}$ factors as
$\rho_{\tilde{M}} = \rho_M \circ p_*$.
The gate axis $\mathbf{n}(\tilde\gamma) = \mathbf{n}(p_*\tilde\gamma)$
depends only on the image in $\pi_1(M)$.
Since the Borel construction depends only on the pairwise dot
products $\mathbf{n}_i \cdot \mathbf{n}_j$, and these are
determined by the holonomy in $\pi_1(M)$, the resulting mixing
matrix is invariant under $p$.
The $\mathbb{Z}/5$-preservation ensures the flat $\mathrm{U}(1)$
CP phase $e^{2\pi i/5}$ is preserved throughout.
\end{proof}

\subsection{Numerical verification}

\begin{center}
\begin{tabular}{lccc}
\toprule
Level & Best triple & Fitness & Max element deviation \\
\midrule
Base (idx 1) &
  \texttt{\{aBaB, baab, baaB\}} & $0.01897$ & --- \\
Deg-2 (idx 39) &
  \texttt{\{bbb, aBBB, bAAb\}} & $0.01897$ & $< 10^{-4}$ \\
Deg-3 (idx 238) &
  \texttt{\{AAb, aaaa, aBAb\}} & $0.01897$ & $< 10^{-4}$ \\
\bottomrule
\end{tabular}
\end{center}

All three levels produce identical $|U_\mathrm{geom}|$ to four
decimal places, with maximum element deviation below $10^{-4}$.
The optimal word triples differ between levels, but produce the
same gate configuration up to numerical precision.

\section{Homology class no-go theorem}
\label{sec:nogo}

\subsection{Flat $\mathrm{U}(1)$ holonomy and homology classes}

For a compact hyperbolic 3-manifold with $H_1(M) \cong \mathbb{Z}/p$,
the flat $\mathrm{U}(1)$ connections are classified by
$\mathrm{Hom}(\mathbb{Z}/p, \mathrm{U}(1))$. The character
$\chi_k$ assigns to each word $\gamma \in \pi_1(M)$ the holonomy
\[
  \chi_k(\gamma) = e^{2\pi i k [\gamma]/p},
\]
where $[\gamma] \in \mathbb{Z}/p$ is the homology class of $\gamma$.

\subsection{The no-go theorem}

\begin{theorem}[Homology class no-go]
\label{thm:nogo}
Let $\{\gamma_1, \gamma_2, \gamma_3\} \subset \pi_1(M)$ be a
word triple with $H_1(M) = \mathbb{Z}/p$. If
$[\gamma_i] = [\gamma_j]$ for some $i \neq j$, then the
Jarlskog CP-violation invariant $J = 0$ identically for any
mixing matrix constructed from this triple via a unitary
orthogonalization of the overlap matrix.
\end{theorem}

\begin{proof}
If $[\gamma_i] = [\gamma_j]$, then
$\chi_k(\gamma_i) = \chi_k(\gamma_j)$ for all $k \in \mathbb{Z}/p$.
In the $3\times 3$ overlap matrix $O_{mn} =
\exp(-\theta_{mn}^2/2\sigma^2)$, the homology-weighted columns
$i$ and $j$ are identical: $O_{\cdot i} = O_{\cdot j}$.
An overlap matrix with two identical columns has rank $\leq 2$.
In the QR decomposition $O = QR$, the matrix $Q$ is unitary
and $R$ is upper triangular. When $O$ has rank $\leq 2$, one
diagonal entry of $R$ vanishes, so $\det R = 0$ and consequently
one column of $Q$ is determined by the other two up to a real
phase. Explicitly, $Q_{m3}$ is a real linear combination of
$Q_{m1}$ and $Q_{m2}$, so all $2\times 2$ minors of $Q$ involving
column~3 are real. The Jarlskog invariant
$J = \mathrm{Im}(U_{11}U_{22}U_{12}^*U_{21}^*)$
is a sum of products of such minors, hence $J = 0$.
\end{proof}

\subsection{Application to the quark sector}

For $m006 = \texttt{OrientableClosedCensus}[43]$ with $H_1 =
\mathbb{Z}/5$, the optimal CKM word triple has homology classes:

\begin{center}
\begin{tabular}{lc}
\toprule
Word & $[\cdot] \in H_1(m006) = \mathbb{Z}/5$ \\
\midrule
\texttt{aaB} & $3$ \\
\texttt{AbA} & $2$ \\
\texttt{AAb} & $2$ \\
\bottomrule
\end{tabular}
\end{center}

The collision $[\texttt{AbA}] = [\texttt{AAb}] = 2$ forces $J = 0$
by Theorem~\ref{thm:nogo}. This is confirmed by the identical
traces: $\mathrm{Tr}(\rho(\texttt{AbA})) =
\mathrm{Tr}(\rho(\texttt{AAb})) = -0.1231 + 2.0108i$, placing
both words in the same conjugacy class of $\mathrm{SU}(2)$.
Verified numerically for $10^3$ random same-trace triples
($J = 0$ to machine precision in all cases).

For the lepton sector ($m003$), the optimal PMNS triple has three
distinct homology classes; no column degeneracy occurs and
$J \neq 0$ is permitted.

\begin{corollary}[Quark-lepton CP asymmetry]
\label{cor:asymmetry}
The CP suppression in the quark sector ($J_\mathrm{CKM} \approx
3\times 10^{-5}$) relative to the lepton sector ($J_\mathrm{PMNS}
\approx 0.031$) is a consequence of Theorem~\ref{thm:nogo}
applied to the optimal gate triples of $m006$ and $m003$
respectively. It is a topological obstruction, not a numerical
coincidence.
\end{corollary}

\section{Discussion}
\label{sec:discussion}

\subsection{Relation to topological quantum computing}

In topological quantum computing~\cite{Kitaev2003,Nayak2008},
gates are realized by braiding non-Abelian anyons, with the gate
determined by the braid group representation. The protection
arises because the gate depends only on the topological class of
the braid, not its geometric details.

The covering tower invariance (Theorem~\ref{thm:invariance})
provides an analogous form of protection: the PMNS gate
configuration depends only on the topological class of the
covering, not on which specific words realize the optimal triple
at each level. Different words at each level produce the same
gate configuration.

This is a different form of protection than braid-based
topological quantum computing: it is \emph{discrete} (invariant
under covering maps, not continuous deformations) rather than
continuous (invariant under isotopy of braids). Whether this
discrete protection has implications for physically realizable
gate operations is an open problem.

\subsection{The no-go theorem as a resource theory}

Theorem~\ref{thm:nogo} establishes that homology class collision
is an obstruction to CP-violating quantum operations in this
framework. This can be reframed as a resource theory: the
\emph{homology class diversity} of a gate triple is a resource
for producing non-zero $J$. Triples with all distinct homology
classes are the resourceful states; triples with collision are
the free states.

This provides a new connection between algebraic topology
(the abelianization $\pi_1(M) \to H_1(M)$) and quantum
information theory (the Jarlskog invariant as a measure of
CP-violating entanglement capacity).

\subsection{What is not claimed}

We do not claim that the PMNS mixing matrix arises from a
physically realizable quantum circuit acting on a qubit. The
construction is mathematical: it produces a unitary matrix from
geometric data. Whether any physical qubit system realizes these
gates is an open problem requiring identification of a suitable
physical platform (e.g., photonic systems with geometric phase,
topological insulators, or Chern-Simons anyons).

\section{Conclusion}
\label{sec:conclusion}

We have established that the holonomy of compact hyperbolic
3-manifolds defines natural qubit rotation gates on the Bloch
sphere, and that specific gate configurations arising from the
lepton and quark sector manifolds of the hyperbolic flavor
geometry framework satisfy two structural theorems.

The covering tower invariance theorem shows that the PMNS gate
configuration is a topological invariant of the
$\mathbb{Z}/5$-preserving covering tower of $m003$: identical
mixing matrices are obtained from three different manifolds
related by covering maps, with maximum element deviation below
$10^{-4}$.

The homology class no-go theorem provides a topological
obstruction to CP-violating quantum operations: gate triples
with homology class collision have $J = 0$ identically. The
quark sector realizes this obstruction exactly, explaining the
observed CP suppression as a topological consequence rather
than a numerical coincidence.

These results connect hyperbolic 3-manifold topology, qubit
gate theory, and Standard Model flavor physics through the
shared mathematical structure of $\mathrm{PSL}(2,\mathbb{C})$
holonomy representations.

\section*{Data Availability}
All scripts are available at
\url{https://github.com/drmlgentry/hyperbolic-flavor-scan}.
The scripts \texttt{qubit\_gates\_v2.py} and
\texttt{borel\_covers\_tower.py} reproduce all numerical
results.

\section*{Conflict of Interest}
The author declares no conflicts of interest.

\begin{thebibliography}{99}

\bibitem{Kitaev2003}
  A.~Y.~Kitaev,
  ``Fault-tolerant quantum computation by anyons,''
  Ann.\ Phys.\ \textbf{303}, 2--30 (2003).

\bibitem{Nayak2008}
  C.~Nayak, S.~H.~Simon, A.~Stern, M.~Freedman, and S.~Das~Sarma,
  ``Non-Abelian anyons and topological quantum computation,''
  Rev.\ Mod.\ Phys.\ \textbf{80}, 1083--1159 (2008).

\bibitem{Gentry:PMNS}
  M.~L.~Gentry,
  ``Lepton Mixing from the Borel Structure of Hyperbolic Holonomy,''
  SSRN preprint, doi:10.2139/ssrn.6583553 (2026).

\bibitem{Gentry:CKM}
  M.~L.~Gentry,
  ``Quark Mixing from Hyperbolic Geometry,''
  SSRN preprint, doi:10.2139/ssrn.6583550 (2026).

\bibitem{Gentry:CP}
  M.~L.~Gentry,
  ``CP Violation from A-Factor Twist Angles of Hyperbolic Holonomy,''
  SSRN preprint, doi:10.2139/ssrn.6583551 (2026).

\bibitem{SnapPy}
  M.~Culler, N.~M.~Dunfield, M.~Goerner, and J.~R.~Weeks,
  ``SnapPy, a computer program for studying the geometry and
  topology of 3-manifolds,''
  \url{http://snappy.computop.org}.

\end{thebibliography}

\end{document}
